\documentclass[11pt,a4paper]{article}

% ============================================================
% PACKAGES
% ============================================================
\usepackage[utf8]{inputenc}
\usepackage[T1]{fontenc}
\usepackage[english]{babel}
\usepackage{amsmath,amssymb,amsthm,mathtools}
\usepackage{mathrsfs}
\usepackage{tikz-cd}
\usepackage{enumitem}
\usepackage{hyperref}
\usepackage{geometry}
\usepackage{microtype}
\usepackage{booktabs}
\usepackage{longtable}
\usepackage{array}
\usepackage{listings}
\usepackage{xcolor}

\geometry{margin=1in}

\lstset{
  basicstyle=\ttfamily\small,
  breaklines=true,
  frame=single,
  numbers=left,
  numberstyle=\tiny,
  keywordstyle=\color{blue!70!black},
  commentstyle=\color{green!50!black},
  stringstyle=\color{red!60!black},
}

% ============================================================
% THEOREM ENVIRONMENTS
% ============================================================
\theoremstyle{definition}
\newtheorem{definition}{Definition}[section]
\newtheorem{cycle}{Cycle}[section]
\newtheorem{principle}[definition]{Principle}
\newtheorem{example}[definition]{Example}
\newtheorem{remark}[definition]{Remark}
\newtheorem{invariant}[definition]{Invariant}
\newtheorem{condition}[definition]{Condition}
\newtheorem{axiom}[definition]{Axiom}

\theoremstyle{plain}
\newtheorem{proposition}[definition]{Proposition}
\newtheorem{theorem}[definition]{Theorem}
\newtheorem{lemma}[definition]{Lemma}
\newtheorem{corollary}[definition]{Corollary}

% ============================================================
% MACROS (inherited from previous papers)
% ============================================================
\newcommand{\Lam}{\Lambda}
\newcommand{\Sp}{\mathrm{Sp}}
\newcommand{\Ge}{\mathrm{Ge}}
\newcommand{\To}{\mathrm{To}}
\newcommand{\Al}{\mathrm{Al}}
\newcommand{\Nu}{\mathrm{Nu}}
\newcommand{\Fo}{\mathrm{Fo}}
\newcommand{\Cl}{\mathrm{Cl}}
\newcommand{\Corr}{\mathrm{Corr}}
\newcommand{\Inv}{\mathrm{Inv}}
\newcommand{\An}{\mathrm{An}}
\newcommand{\Cd}{\mathrm{Cd}}
\newcommand{\Var}{\mathrm{Var}}
\newcommand{\Ob}{\mathfrak{O}}
\newcommand{\Xcal}{\mathcal{X}}
\newcommand{\Ccal}{\mathcal{C}}
\newcommand{\Rcal}{\mathcal{R}}
\newcommand{\Mcal}{\mathcal{M}}
\newcommand{\Ncal}{\mathcal{N}}
\newcommand{\Pcal}{\mathcal{P}}
\newcommand{\Hcal}{\mathcal{H}}
\newcommand{\Acal}{\mathcal{A}}
\newcommand{\Dcal}{\mathcal{D}}
\newcommand{\Tcal}{\mathcal{T}}
\newcommand{\Vcal}{\mathcal{V}}
\newcommand{\Ical}{\mathcal{I}}
\newcommand{\Ecal}{\mathcal{E}}
\newcommand{\Ocal}{\mathcal{O}}
\newcommand{\Fcal}{\mathcal{F}}
\newcommand{\Rfrak}{\mathfrak{R}}
\newcommand{\Cfrak}{\mathfrak{C}}
\newcommand{\philam}{\phi_{\lambda\mu}}
\newcommand{\Vmat}{\mathbf{V}}
\newcommand{\Jset}{\mathcal{J}}
\newcommand{\dd}{\,\mathrm{d}}
\newcommand{\op}[1]{\operatorname{#1}}
\newcommand{\spec}{\sigma}
\newcommand{\Norm}[1]{\left\lVert #1 \right\rVert}
\newcommand{\Abs}[1]{\left\lvert #1 \right\rvert}
\newcommand{\ip}[2]{\left\langle #1, #2 \right\rangle}
\newcommand{\R}{\mathbb{R}}
\newcommand{\C}{\mathbb{C}}
\newcommand{\N}{\mathbb{N}}
\newcommand{\Z}{\mathbb{Z}}
\newcommand{\Q}{\mathbb{Q}}
\newcommand{\Pmark}{\mathrm{P}}
\newcommand{\Nmark}{\mathrm{N}}
\newcommand{\Omark}{\mathrm{O}}
\newcommand{\Fmark}{\mathrm{F}}
\newcommand{\STPASS}{\textsc{Pass}}
\newcommand{\STFAIL}{\textsc{Fail}}
\newcommand{\STUNRES}{\textsc{Unresolved}}
\newcommand{\STCLOSED}{\textsc{Closed}}
\newcommand{\STOPEN}{\textsc{Open}}
\newcommand{\STOBSTR}{\textsc{Obstructed}}
\newcommand{\STIMPLIED}{\textsc{Code-Implied}}
\newcommand{\STVERIFIED}{\textsc{Code-Verified}}

% ============================================================
% TITLE
% ============================================================
\title{%
  \textbf{From Code to Structure:}\\[4pt]
  \large An Operational Protocol for Reverse-Engineering\\
  Mathematical Objects from Computational Artifacts
}
\author{}
\date{}

\begin{document}
\maketitle

\begin{abstract}
\noindent
We present an operational protocol for analysing mathematical
problems whose point of entry is executable code rather than a
formal statement. The code is treated as a computational
\emph{artifact}, not as a mathematical object: it produces outputs
but does not justify the choices, hypotheses, or approximations
that make those outputs meaningful. The protocol reverse-engineers
the mathematical object that the code approximates and certifies
the approximation. It is organised as eight mandatory cycles
(C0--C7) and terminates in exactly one of three states:
\textsc{Closed}, \textsc{Open}, or \textsc{Obstructed}. The output
of the protocol is not a paper but a \emph{certification
document}: a structured record comprising the code, the inferred
model, the reconstructed continuous object, the analysis by layer,
an independent re-discretisation, a cross-audit matrix, and a
provenance trace. We formalise six validity conditions (V1--V6)
that must hold for the certificate to be complete, and we argue
that the code becomes \emph{redundant} precisely when each of its
outputs is predictable from the reconstructed structure without
execution. The protocol is explicitly positioned with respect to
four companion papers: a framework paper on multirepresentational
constraining, a paper on structural constraining and spectral
emergence, a paper on the limits of constraining, and a position
paper on structural auditability.
\end{abstract}

\tableofcontents

% ============================================================
\part{Introduction and Positioning}
% ============================================================

\section{Motivation: the Code-First Problem}
\label{sec:motivation}

\subsection{The Situation}

A researcher receives a piece of code. The code runs, produces
numbers, plots, tables. It ``works''. The question is not whether
the code produces outputs; it evidently does. The question is what
those outputs \emph{mean}.

The code contains, simultaneously:
\begin{itemize}
  \item mathematical content (an equation, an operator, boundary
        conditions);
  \item computational content (a scheme, a grid, a stopping
        criterion, a floating-point arithmetic).
\end{itemize}
It fuses the two. A reader cannot, by inspecting the code alone,
separate the mathematical object from its approximation. Two
readers may interpret the same code as solving two different
equations, with the same numerical output.

\begin{definition}[Computational artifact]\label{def:artifact}
A \emph{computational artifact} is a pair
\[
  \Acal = (\mathrm{src}, \mathrm{out})
\]
where $\mathrm{src}$ is executable source code and $\mathrm{out}$
is a collection of observable outputs (numbers, plots, tables,
error logs) produced by executing $\mathrm{src}$.
\end{definition}

\begin{remark}
The artifact is \emph{not} a mathematical object. It is a witness
to the existence of a numerical procedure. A witness is not a
proof.
\end{remark}

\subsection{The Question the Protocol Answers}

Given an artifact $\Acal$, the protocol answers:
\begin{enumerate}[label=(\roman*)]
  \item What mathematical object does $\Acal$ approximate?
  \item Is that object well-defined?
  \item Does the approximation converge to the object?
  \item Is the object consistent with its structural invariants?
  \item Can the code be \emph{reconstructed} from the object, so
        that its outputs are predictable without execution?
\end{enumerate}

\subsection{The Inversion of Direction}

Classical mathematical work proceeds from problem to code:
\[
  \text{problem} \;\longrightarrow\; \text{representations}
  \;\longrightarrow\; \text{code} \;\longrightarrow\;
  \text{verification}.
\]
The protocol proceeds in the opposite direction:
\[
  \text{code} \;\longrightarrow\; \text{inferred model}
  \;\longrightarrow\; \text{continuous object}
  \;\longrightarrow\; \text{structure} \;\longrightarrow\;
  \text{validation against the code}.
\]
The code becomes the \emph{object of study}, not the tool.

\subsection{The Central Criterion}

The protocol is successful when the code becomes
\emph{redundant}:

\begin{principle}[Code redundancy]\label{pr:redundancy}
The protocol has completed its analysis of an artifact $\Acal$
when each output in $\mathrm{out}$ is predictable from the
reconstructed continuous object and its structural layers, without
executing $\mathrm{src}$.
\end{principle}

This is the operational criterion of closure. It replaces the
vaguer criterion ``the analysis is complete'' with a checkable
condition: recompute the outputs from the structure, and compare.

% ------------------------------------------------------------
\section{Relationship to Companion Papers}
\label{sec:relationship}
% ------------------------------------------------------------

\subsection{The Five Papers}

This paper belongs to a program developed across five previous
contributions:

\begin{center}
\begin{tabular}{l l}
  \toprule
  Reference & Focus \\
  \midrule
  \cite{framework} & Framework: master object, layers, transitions, \\
                   & reconstruction, validation matrix. \\
  \cite{paperI}    & Structural constraining, spectral emergence, \\
                   & discrepancy trichotomy; four worked toys. \\
  \cite{paperII}   & Separation criterion, three-type rupture, \\
                   & cross-coherence, dynamic expansion. \\
  \cite{paperIII}  & Structural auditability, wear, certification \\
                   & trace, six-level output taxonomy. \\
  \textbf{This paper} & Operational protocol for code-first \\
                   & analysis; certification document. \\
  \bottomrule
\end{tabular}
\end{center}

\subsection{What is Inherited}

From \cite{framework}:
\begin{itemize}
  \item the master object $\Ob$ and the layer set $\Lam$;
  \item the transition morphisms $\philam$;
  \item the reconstruction operators $\rho_\Cl, \rho_\Inv, \rho_\An$;
  \item the validation matrix $\Vmat$ and the epistemic states
        $\{\Pmark, \Nmark, \Omark, \Fmark\}$;
  \item the discrete family $\Ob^{(N)}$ and the error
        $\varepsilon_N$.
\end{itemize}

From \cite{paperI}:
\begin{itemize}
  \item structural constraining: $\Ccal_0 \supseteq \Ccal_1
        \supseteq \cdots$;
  \item spectral emergence via Fr\'echet linearisation;
  \item discrepancy trichotomy (discretisation / extension /
        inconsistency).
\end{itemize}

From \cite{paperII}:
\begin{itemize}
  \item separation criterion $\Phimap(X) = \Phimap(Y)
        \Rightarrow X \simadm Y$;
  \item three-type rupture classification (Type I, II, III);
  \item coherence subspace $\Rfrak_{\rm coh} \subset \prod_i
        \Rcal_i$;
  \item dynamic expansion $\Rfrak_{n+1} = \Rfrak_n \cup
        \{R_{\rm new}\}$.
\end{itemize}

From \cite{paperIII}:
\begin{itemize}
  \item the six-level output taxonomy;
  \item the certification trace $\Tcal(X) = (D, H, C, R, E, V, A)$;
  \item the principles of auditability, hostile representation,
        and trace completeness.
\end{itemize}

\subsection{What is New}

\begin{enumerate}[label=(\roman*)]
  \item The \textbf{Code layer} $R_{\Cd}$ as a formal input.
  \item The \textbf{inference morphisms}
        $\iota_{\Cd \to \Ge}$, $\iota_{\Cd \to \Nu}$.
  \item The \textbf{eight-cycle protocol} C0--C7.
  \item The \textbf{six validity conditions} V1--V6.
  \item The \textbf{certification document} as the formal output.
  \item The \textbf{code redundancy criterion} (Principle~\ref{pr:redundancy}).
\end{enumerate}

% ------------------------------------------------------------
\section{Contributions and Organization}
\label{sec:contributions}
% ------------------------------------------------------------

\subsection{Contributions}

\begin{enumerate}[label=(\roman*)]
  \item A formal definition of the code layer and its inference
        morphisms (\S\ref{sec:foundations-code}).
  \item An eight-cycle protocol for code-first analysis
        (\S\ref{sec:protocol}).
  \item Six validity conditions V1--V6 defining complete
        certification (\S\ref{sec:validity}).
  \item A specification of the certification document
        (\S\ref{sec:certificate}).
  \item Two worked examples: a \textsc{Closed} case and an
        \textsc{Obstructed} case (\S\ref{sec:cases}).
  \item A backend architecture for implementation
        (\S\ref{sec:backend}).
\end{enumerate}

\subsection{Organization}

Part~\ref{part:intro} positions the protocol. Part~\ref{part:foundations}
recalls the framework and introduces the code layer. Part~\ref{part:protocol}
specifies the eight cycles. Part~\ref{part:validity} formalises the
validity conditions. Part~\ref{part:certificate} specifies the
certificate. Part~\ref{part:cases} presents two worked cases.
Part~\ref{part:backend} sketches the implementation. Part~\ref{part:closing}
closes.

% ============================================================
\part{Foundations}
\label{part:foundations}
% ============================================================

\section{Recap of the Multirepresentational Framework}
\label{sec:foundations-recap}

We recall the essential structure of \cite{framework} without
reproducing its full content.

\subsection{The Master Object and Layers}

\begin{definition}[Master object, {\cite[Definition 1.1]{framework}}]
A \emph{master object} is a tuple
\[
  \Ob = (\Ocal_\lambda)_{\lambda \in \Lam},
\]
where
\[
  \Lam = \{\Sp, \Ge, \To, \Al, \Nu, \Fo, \Cl, \Corr, \An\}
\]
and each $\Ocal_\lambda$ is a layer datum.
\end{definition}

\subsection{Transition and Reconstruction}

Transition morphisms $\philam : \Ocal_\lambda \to \Ocal_\mu$ for
$\lambda \prec \mu$ carry structure between layers. The
construction pipeline is
\[
  \Ocal_\Sp \to \Ocal_\Ge \to \Ocal_\To \to \Ocal_\Al
  \to \Ocal_\Nu \to \Ocal_\Fo .
\]
The reconstruction operators
\[
  \rho_\Cl : \Ocal_\Fo \to \Ecal_\Cl, \quad
  \rho_\Inv : \Ecal_\Cl \to \Ecal_\Corr, \quad
  \rho_\An : \Ecal_\Corr \to \Acal
\]
return from the formal layer to the classical and corrected
models.

\subsection{The Validation Matrix}

\begin{definition}[Validation matrix, {\cite[Definition 1.6]{framework}}]
The matrix
\[
  \Vmat = (v_{\lambda j})_{\lambda \in \Lam, \, j \in \Jset},
  \qquad
  v_{\lambda j} \in \{\Pmark, \Nmark, \Omark, \Fmark\},
\]
with
\[
  \Jset = \{\mathrm{Ex}, \mathrm{Co}, \mathrm{In}, \mathrm{Nu},
          \mathrm{St}, \mathrm{Cv}, \mathrm{Re}, \mathrm{Rev}\}
\]
records the epistemic status of each criterion at each layer.
\end{definition}

\subsection{Constraining}

From \cite{paperI}, the constrained model spaces form a decreasing
chain:
\[
  \Ccal_0 \supseteq \Ccal_\Sp \supseteq \Ccal_\Ge \supseteq \cdots
  \supseteq \Ccal_\An .
\]
The methodology's output is the intersection
\[
  \Ccal_\ast = \bigcap_{\lambda \in \Lam} \Ccal_\lambda .
\]

\subsection{Separation}

From \cite{paperII}, the representation map
\[
  \Phimap : \Xcal \to \prod_{\lambda \in \Lam} \Rcal_\lambda,
  \qquad
  \Phimap(X) = (R_\lambda(X))_{\lambda \in \Lam}
\]
must satisfy
\[
  \Phimap(X) = \Phimap(Y) \;\Longrightarrow\; X \simadm Y
\]
for the methodology to close. Failure of this condition defines a
\emph{rupture} of Type I, II, or III.

\subsection{Auditability}

From \cite{paperIII}, every construction should be accompanied by
a certification trace
\[
  \Tcal(X) = (D, H, C, R, E, V, A),
\]
recording definitions, hypotheses, constructions, representations,
estimates, verifications, and audits.

% ------------------------------------------------------------
\section{The Code Layer}
\label{sec:foundations-code}
% ------------------------------------------------------------

\subsection{Formal Extension of the Layer Set}

\begin{definition}[Extended layer set]\label{def:extended-layers}
The \emph{extended layer set} is
\[
  \Lam^+ = \Lam \cup \{\Cd\},
\]
where $\Cd$ denotes the \emph{code layer}.
\end{definition}

\begin{definition}[Code datum]\label{def:code-datum}
A \emph{code datum} is a tuple
\[
  \Ocal_\Cd = (\mathrm{src}, \mathrm{out}, \mathrm{env}),
\]
where:
\begin{itemize}
  \item $\mathrm{src}$ is the source text;
  \item $\mathrm{out}$ is the set of observable outputs;
  \item $\mathrm{env}$ is the execution environment (language,
        libraries, precision, hardware assumptions).
\end{itemize}
\end{definition}

\subsection{Inference Morphisms}

\begin{definition}[Inference morphism]\label{def:inference}
An \emph{inference morphism} is a map
\[
  \iota_{\Cd \to \lambda} : \Ocal_\Cd \to \Ocal_\lambda
\]
from the code layer to a structural layer $\lambda$, together with
a \emph{justification} $J_{\Cd \to \lambda}$ recording the
assumptions under which the inference is valid.
\end{definition}

Two inference morphisms are primary:

\begin{itemize}
  \item $\iota_{\Cd \to \Ge} : \Ocal_\Cd \to \Ocal_\Ge$: from code
        to the inferred geometric datum (domain, dimension,
        boundary structure).
  \item $\iota_{\Cd \to \Nu} : \Ocal_\Cd \to \Ocal_\Nu$: from code
        to the inferred numerical datum (scheme, grid, tolerance,
        convergence order).
\end{itemize}

\begin{axiom}[Inference non-uniqueness]\label{ax:non-unique}
In general, inference morphisms are not unique: multiple structural
objects may be compatible with a single code datum.
\end{axiom}

\begin{remark}
Axiom~\ref{ax:non-unique} is the formal statement of the
code-first problem: the code \emph{underdetermines} the object it
approximates. The protocol's role is to reduce this
underdetermination via additional representations.
\end{remark}

\subsection{The Inferred Model}

\begin{definition}[Inferred model]\label{def:inferred-model}
Given an inference morphism $\iota_{\Cd \to \Ge}$, the
\emph{inferred model} is
\[
  \Mcal_{\rm imp} = \iota_{\Cd \to \Ge}(\Ocal_\Cd) \in \Ocal_\Ge,
\]
together with a \emph{hypothesis record}
$\Hcal_{\rm imp} = (J_{\Cd \to \Ge}, \text{assumptions})$.
\end{definition}

The inferred model is a \emph{candidate} for the continuous object.
It becomes a genuine identification only after the protocol's
validation cycles.

% ------------------------------------------------------------
\section{Extended Epistemic States}
\label{sec:epistemic}
% ------------------------------------------------------------

The epistemic states of \cite{framework} are extended to include
two code-specific states.

\begin{definition}[Extended epistemic states]\label{def:epistemic-states}
The \emph{extended epistemic state set} is
\[
  \mathrm{Ep} = \{\Pmark, \Nmark, \Omark, \Fmark,
                  \STIMPLIED, \STVERIFIED\},
\]
where:
\begin{itemize}
  \item $\Pmark$: proved mathematically;
  \item $\Nmark$: numerically supported;
  \item $\Omark$: open;
  \item $\Fmark$: failed;
  \item $\STIMPLIED$: implied by the code, without justification;
  \item $\STVERIFIED$: verified against the code via an independent
        re-discretisation.
\end{itemize}
\end{definition}

\begin{principle}[No silent promotion]\label{pr:no-silent}
A statement marked $\STIMPLIED$ may be promoted to $\Pmark$,
$\Nmark$, or $\STVERIFIED$ only by exhibiting the explicit
mathematical or computational argument that establishes the
promotion. Silent promotion is a protocol violation.
\end{principle}

% ============================================================
\part{The Protocol}
\label{part:protocol}
% ============================================================

\section{Overview}
\label{sec:protocol-overview}

\subsection{Structure}

The protocol consists of eight mandatory cycles
\[
  \mathrm{C0} \to \mathrm{C1} \to \cdots \to \mathrm{C7},
\]
each with a fixed signature: input, operation, output, success
criterion, failure modes, invariant, and transition. Cycles cannot
be reordered. Sub-operations within a cycle may be parallelised.

\subsection{Terminal States}

The protocol terminates in exactly one of three states:

\begin{itemize}
  \item \STCLOSED: the residue reduces to a single equivalence
        class and the reconstruction reproduces the code.
  \item \STOPEN: the residue contains multiple classes but no
        obstruction is identified.
  \item \STOBSTR: a Type~III rupture (\cite{paperII},
        Definition~3.6) is identified.
\end{itemize}

\subsection{Global Invariants}

\begin{invariant}[Provenance]\label{inv:provenance}
Every result produced by any cycle carries a provenance record
identifying the cycle, its input, its operation, and the
assumptions used. This is the cycle-level realisation of the
certification trace of \cite{paperIII}, Definition~6.1.
\end{invariant}

\begin{invariant}[Epistemic separation]\label{inv:epistemic}
Every result is tagged with exactly one state in $\mathrm{Ep}$
(Definition~\ref{def:epistemic-states}). $\Pmark$ and $\Nmark$ are
never conflated.
\end{invariant}

\begin{invariant}[Strict residue decrease]\label{inv:residue}
Each iteration of the loop
$\mathrm{C1} \to \mathrm{C2} \to \mathrm{C3} \to \mathrm{C4}
\to \mathrm{C5} \to \mathrm{C6}$ strictly decreases the residue:
\[
  \Ccal^{(n+1)} \subsetneq \Ccal^{(n)}.
\]
\end{invariant}

% ------------------------------------------------------------
\section{Cycle C0: Code Reading}
\label{sec:c0}
% ------------------------------------------------------------

\begin{cycle}[Code reading]\label{cy:c0}
\textbf{Input.} A computational artifact $\Acal = (\mathrm{src},
\mathrm{out})$ (Definition~\ref{def:artifact}).

\medskip
\textbf{Operation.}
\begin{enumerate}[label=\arabic*.]
  \item Parse $\mathrm{src}$ to extract: variables, operations,
        loops, data structures, constants, tolerances.
  \item Identify the \emph{governing hypothesis}: what equation,
        operator, or functional is the code approximating.
  \item Enumerate the \emph{implicit choices}: discretisation
        scheme, grid, boundary treatment, stopping criterion,
        floating-point precision.
  \item Record $\mathrm{out}$ and its structure (which outputs are
        scalar, which are arrays, which are time series, etc.).
\end{enumerate}

\medskip
\textbf{Output.} A structured hypothesis
\[
  \Hcal_0 = (\text{governing hypothesis}, \text{implicit choices},
            \text{output structure})
\]
with a provenance record.

\medskip
\textbf{Success criterion.} $\Hcal_0$ is complete and unambiguous:
every field is explicit.

\medskip
\textbf{Failure modes.}
\begin{itemize}
  \item Ambiguity in the governing hypothesis: return to input.
  \item Obfuscated code (no semantic structure recoverable): flag
        as \STOPEN.
  \item Multiple incompatible readings: emit all candidates, defer
        disambiguation to C1.
\end{itemize}

\medskip
\textbf{Invariant.} The governing hypothesis must be stated as a
mathematical relation, not as a procedural description.

\medskip
\textbf{Transition.} $\to$ C1.
\end{cycle}

\begin{remark}
Cycle C0 is the analogue, for code-first analysis, of
\cite[Step P0]{paperI}. In \cite{paperI}, P0 formalised an
informal problem; here, C0 formalises an \emph{artifact}. The
direction of the operation is inverted.
\end{remark}

% ------------------------------------------------------------
\section{Cycle C1: Model Inference}
\label{sec:c1}
% ------------------------------------------------------------

\begin{cycle}[Model inference]\label{cy:c1}
\textbf{Input.} The hypothesis $\Hcal_0$.

\medskip
\textbf{Operation.} Construct the inferred model
\[
  \Mcal_{\rm imp} = \iota_{\Cd \to \Ge}(\Ocal_\Cd)
\]
as a mathematical object: domain, governing equation, boundary
conditions, functional spaces, and hypotheses.

For each implicit choice in $\Hcal_0$, produce a corresponding
structural declaration:
\begin{itemize}
  \item discretisation scheme $\to$ order of consistency;
  \item grid $\to$ regularity of the domain;
  \item boundary treatment $\to$ boundary conditions;
  \item stopping criterion $\to$ tolerance parameter.
\end{itemize}

\medskip
\textbf{Output.} The inferred model $\Mcal_{\rm imp}$ together with
the hypothesis record $\Hcal_{\rm imp}$ (Definition~\ref{def:inferred-model}).

\medskip
\textbf{Success criterion.} $\Mcal_{\rm imp}$ is a well-defined
mathematical object: domain, equation, and conditions are
unambiguous.

\medskip
\textbf{Failure modes.}
\begin{itemize}
  \item $\Mcal_{\rm imp}$ not well-defined: return to C0 with the
        specific ambiguity.
  \item Multiple well-defined candidates: emit each as a separate
        branch and defer to C2.
\end{itemize}

\medskip
\textbf{Invariant.} Every field of $\Mcal_{\rm imp}$ must be
declared as either inferred with justification or assumed without
justification. The distinction is recorded via Principle~\ref{pr:no-silent}.
\end{invariant}

\medskip
\textbf{Transition.} $\to$ C2.
\end{cycle}

% ------------------------------------------------------------
\section{Cycle C2: Classical Theory Identification}
\label{sec:c2}
% ------------------------------------------------------------

\begin{cycle}[Classical theory identification]\label{cy:c2}
\textbf{Input.} The inferred model $\Mcal_{\rm imp}$.

\medskip
\textbf{Operation.} Search the mathematical literature for the
object $\Mcal_{\rm imp}$. Determine:
\begin{enumerate}[label=\arabic*.]
  \item Whether $\Mcal_{\rm imp}$ is a known object (identified by
        standard notation and reference).
  \item What is already known about $\Mcal_{\rm imp}$: spectral
        theory, existence, uniqueness, regularity, asymptotic
        behaviour.
  \item Whether the code's hypotheses are compatible with the
        classical treatment.
\end{enumerate}

\medskip
\textbf{Output.} A classification
\[
  \text{status}(\Mcal_{\rm imp}) \in
  \{\textsc{Known}, \textsc{Variant}, \textsc{New}\},
\]
with references for the \textsc{Known} and \textsc{Variant} cases.

\medskip
\textbf{Success criterion.} The status is determined, and for
\textsc{Known} cases the reference is identified.

\medskip
\textbf{Failure modes.}
\begin{itemize}
  \item The object cannot be matched to any classical theory: mark
        as \textsc{New}, proceed to C3.
  \item The object matches multiple incompatible theories: emit
        both, defer disambiguation to C3.
\end{itemize}

\medskip
\textbf{Invariant.} If the status is \textsc{Known}, the classical
properties of the object (existence, uniqueness, regularity,
spectral theory) are inherited automatically and need not be
reproved.

\medskip
\textbf{Transition.} $\to$ C3.
\end{cycle}

\begin{remark}
Cycle C2 corresponds, in the direct direction of \cite{framework},
to the identification of the classical model via $\rho_\Cl$. Here
it is performed \emph{before} the structural layers are built, as a
search operation.
\end{remark}

% ------------------------------------------------------------
\section{Cycle C3: Structural Layer Construction}
\label{sec:c3}
% ------------------------------------------------------------

\begin{cycle}[Structural layer construction]\label{cy:c3}
\textbf{Input.} The inferred model $\Mcal_{\rm imp}$.

\medskip
\textbf{Operation.} Build the structural layers for
$\Mcal_{\rm imp}$:
\begin{itemize}
  \item spectral layer $\Ocal_\Sp$: spectrum, eigenfunctions,
        spectral measure;
  \item geometric layer $\Ocal_\Ge$: domain, metric, connection,
        curvature;
  \item topological layer $\Ocal_\To$: homology, homotopy, Euler
        characteristic;
  \item algebraic layer $\Ocal_\Al$: symmetry group, operator
        algebra, factorisation;
  \item variational layer $\Ocal_\Var$: functional, Euler--Lagrange
        operator, energy law.
\end{itemize}
If the spectral layer is not directly available (nonlinear operator),
apply the Fr\'echet linearisation of \cite[Construction 3.6]{paperI}.

\medskip
\textbf{Output.} The family
\[
  \Rcal_n = \{R_\lambda\}_{\lambda \in \Lam_n}
\]
where $\Lam_n \subseteq \Lam^+$ is the set of layers successfully
constructed.

\medskip
\textbf{Success criterion.} At least three layers are non-trivial.
The spectral and variational layers are mandatory when the operator
admits a variational structure.

\medskip
\textbf{Failure modes.}
\begin{itemize}
  \item Fewer than three layers: flag as \STOPEN, return to C1 to
        expand the inferred model.
  \item Spectral emergence required but Fr\'echet linearisation is
        undefined: flag as \STOBSTR.
\end{itemize}

\medskip
\textbf{Invariant.} Every constructed layer must carry a provenance
record and an epistemic state from $\mathrm{Ep}$
(Definition~\ref{def:epistemic-states}).

\medskip
\textbf{Transition.} $\to$ C4.
\end{cycle}

% ------------------------------------------------------------
\section{Cycle C4: Independent Re-discretisation}
\label{sec:c4}
% ------------------------------------------------------------

\begin{cycle}[Independent re-discretisation]\label{cy:c4}
\textbf{Input.} The structural layers $\Rcal_n$.

\medskip
\textbf{Operation.} Build an \emph{independent} numerical scheme
$\Dcal^\ast$ for $\Mcal_{\rm imp}$, using a method that is
\emph{not} the one in the original code. For example:
\begin{itemize}
  \item if the original code uses finite differences, use spectral
        methods or finite elements;
  \item if the original code uses explicit time stepping, use
        implicit or exponential integrators;
  \item if the original code uses a specific grid, use an adaptive
        or unstructured grid.
\end{itemize}
Execute $\Dcal^\ast$ and compare its outputs to $\mathrm{out}$.

\medskip
\textbf{Output.} A comparison record
\[
  \mathrm{Comp} = (\mathrm{out}^\ast, \mathrm{out},
    \text{difference norm}, \text{order of convergence}).
\]

\medskip
\textbf{Success criterion.} The two outputs agree to within the
declared tolerance $\mathcal{T}$ (Definition of $\Pcal_0$ in
\cite[Step P0]{paperI}).

\medskip
\textbf{Failure modes.}
\begin{itemize}
  \item Discrepancy within $\mathcal{T}$: proceed to C5.
  \item Discrepancy outside $\mathcal{T}$: classify the discrepancy
        using the trichotomy of \cite[Theorem 4.3]{paperI}
        (discretisation / extension / inconsistency). Return to C1
        with the classification.
\end{itemize}

\medskip
\textbf{Invariant.} The scheme $\Dcal^\ast$ must be demonstrably
different from the scheme of $\mathrm{src}$. A redundant scheme
(e.g.\ the same method with a different variable name) does not
satisfy the independence requirement.

\medskip
\textbf{Transition.} $\to$ C5.
\end{cycle}

\begin{remark}
Cycle C4 is where the code is put on trial. The original code has
been used to \emph{infer} the model; now an independent scheme tests
whether that inference is correct.
\end{remark}

% ------------------------------------------------------------
\section{Cycle C5: Cross-Audit}
\label{sec:c5}
% ------------------------------------------------------------

\begin{cycle}[Cross-audit]\label{cy:c5}
\textbf{Input.} The family $\Rcal_n$ and the comparison record
$\mathrm{Comp}$.

\medskip
\textbf{Operation.} For every pair $(R_i, R_j)$ of layers, verify
the coherence relation
\[
  (R_i(X), R_j(X)) \in C_{ij}
  \qquad \forall X \in \Ccal^{(n)},
\]
as in \cite[Definition 4.1]{paperII}. Build the audit matrix
\[
  \Acal = (a_{ij}), \qquad a_{ij} \in
  \{\STPASS, \STFAIL, \STUNRES\}.
\]

\medskip
\textbf{Output.} The audit matrix $\Acal$ together with
localisation records for each $\STFAIL$ entry: which pair of
representations, which condition, which evidence.

\medskip
\textbf{Success criterion.} All pairs satisfy
$a_{ij} = \STPASS$.

\medskip
\textbf{Failure modes.}
\begin{itemize}
  \item Some $a_{ij} = \STFAIL$: classify the discrepancy
        (Type I, II, or III) using \cite[Theorem 4.3]{paperI} and
        \cite[\S3]{paperII}. Proceed to C6.
  \item Some $a_{ij} = \STUNRES$: proceed to C6 with an \textsc{Open}
        tag.
\end{itemize}

\medskip
\textbf{Invariant.} The matrix $\Acal$ must be symmetric modulo the
coherence relation: $a_{ij} = a_{ji}$ whenever the coherence is
symmetric.

\medskip
\textbf{Transition.} $\to$ C6.
\end{cycle}

% ------------------------------------------------------------
\section{Cycle C6: Code Reconstruction}
\label{sec:c6}
% ------------------------------------------------------------

\begin{cycle}[Code reconstruction]\label{cy:c6}
\textbf{Input.} The family $\Rcal_n$, the audit matrix $\Acal$, and
the comparison record $\mathrm{Comp}$.

\medskip
\textbf{Operation.} Attempt to reconstruct the original code from
the structural layers:
\[
  \Dcal^\ast_{\rm rec} = \text{reconstruct}(\Rcal_n, \Acal).
\]
Compare $\Dcal^\ast_{\rm rec}$ to the original scheme
$\Dcal_{\rm orig}$ inferred at C1.

Two comparisons are required:
\begin{enumerate}[label=\arabic*.]
  \item \emph{Structural comparison}: does
        $\Dcal^\ast_{\rm rec}$ use the same order of consistency,
        the same regularity hypotheses, the same boundary treatment?
  \item \emph{Output comparison}: are the outputs of
        $\Dcal^\ast_{\rm rec}$ and $\Dcal_{\rm orig}$ identical to
        within $\mathcal{T}$?
\end{enumerate}

\medskip
\textbf{Output.} A reconstruction record
\[
  \mathrm{Rec} = (\Dcal^\ast_{\rm rec},
    \text{structural match}, \text{output match}).
\]

\medskip
\textbf{Success criterion.} Both comparisons succeed. In this
case, the code is \emph{redundant} in the sense of
Principle~\ref{pr:redundancy}.

\medskip
\textbf{Failure modes.}
\begin{itemize}
  \item Structural or output mismatch: classify the discrepancy.
        If Type I or II, return to C3 with the missing layer. If
        Type III, terminate at \STOBSTR.
\end{itemize}

\medskip
\textbf{Invariant.} Reconstruction must be \emph{independent} of
the original code: the reconstruction procedure cannot consult
$\mathrm{src}$.

\medskip
\textbf{Transition.} $\to$ C7 if both comparisons succeed;
otherwise, return to C3 or terminate at \STOBSTR.
\end{cycle}

\begin{remark}
Cycle C6 is the operational test of Principle~\ref{pr:redundancy}.
If the reconstruction succeeds, the code is certified as a
redundant consequence of the structure. If it fails, the failure
localises what part of the structure is not yet captured.
\end{remark}

% ------------------------------------------------------------
\section{Cycle C7: Formalization}
\label{sec:c7}
% ------------------------------------------------------------

\begin{cycle}[Formalization]\label{cy:c7}
\textbf{Input.} The complete trace of cycles C0--C6.

\medskip
\textbf{Operation.} Encode in a proof assistant (Lean~4, Coq,
Isabelle, or Agda) the fragments of the analysis that admit
mechanical verification:
\begin{itemize}
  \item the definition of the inferred model $\Mcal_{\rm imp}$;
  \item the well-posedness of the operator;
  \item the spectral theorem applied to the linearised operator;
  \item the coherence matrix and its verification;
  \item the reconstruction theorem, when expressible.
\end{itemize}

\medskip
\textbf{Output.} A formal certificate consisting of:
\begin{itemize}
  \item source files in the proof assistant's language;
  \item the list of proof obligations discharged;
  \item the list of fragments not formalisable, with reasons.
\end{itemize}

\medskip
\textbf{Success criterion.} At least the following are formalised:
\begin{itemize}
  \item the definition of $\Mcal_{\rm imp}$ and its domain;
  \item the spectral or variational structure of the linearised
        operator;
  \item the coherence relations in the audit matrix $\Acal$.
\end{itemize}

\medskip
\textbf{Failure modes.}
\begin{itemize}
  \item Fragments not formalisable within the current
        infrastructure: recorded as $\Omark$.
  \item Fragments proved false by the assistant: recorded as
        $\Fmark$, return to C1.
\end{itemize}

\medskip
\textbf{Invariant.} The formal certificate must be self-contained:
a third party can replay the proof without consulting the
original code.

\medskip
\textbf{Transition.} $\to$ terminal state.
\end{cycle}

% ============================================================
\part{Validity}
\label{part:validity}
% ============================================================

\section{The Six Validity Conditions}
\label{sec:validity}

The protocol's output is a certificate. The certificate is
\emph{complete} if and only if six conditions hold simultaneously.

\begin{condition}[V1: Reconstruction]\label{cond:v1}
There exists a scheme $\Dcal^\ast$ such that
\[
  \Dcal^\ast(\Mcal_{\rm imp}) \approx \Dcal_{\rm orig}
\]
in both structure and output. This is the output of Cycle C6.
\end{condition}

\begin{condition}[V2: Convergence]\label{cond:v2}
The outputs of $\Dcal^\ast$ converge to those of
$\Mcal_{\rm imp}$ with a \emph{demonstrated} rate:
\[
  \Norm{u_N - u} \leq C N^{-p},
\]
where $C$ and $p$ are proved, not merely estimated. This is the
content of \cite[Theorem 4.6]{paperI} (Lax--Richtmyer
consistency-stability-convergence).
\end{condition}

\begin{condition}[V3: Coherence]\label{cond:v3}
All pairs of layers satisfy
\[
  a_{ij} = \STPASS \qquad \forall (i,j).
\]
This is the output of Cycle C5 and the coherence criterion of
\cite[Definition 4.1]{paperII}.
\end{condition}

\begin{condition}[V4: Independence]\label{cond:v4}
The reconstruction is not circular: the object $\Mcal_{\rm imp}$
is not inferred exclusively from the code. At least one layer
(typically $\Ocal_\Ge$, $\Ocal_\To$, or $\Ocal_\Al$) provides an
independent characterisation.
\end{condition}

\begin{condition}[V5: Auditability]\label{cond:v5}
Every step of the protocol carries a provenance record and an
epistemic state from $\mathrm{Ep}$. The full trace is
reconstructible by a third party. This is the trace-completeness
principle of \cite[Principle 6.4]{paperIII}.
\end{condition}

\begin{condition}[V6: Formalization]\label{cond:v6}
The fragments identified in Cycle C7 have been encoded in a proof
assistant, and the discharged obligations cover the definitions of
the inferred model, the structural layers, and the reconstruction.
\end{condition}

\begin{theorem}[Complete certificate]\label{th:complete}
The certification document is complete if and only if
\[
  V1 \wedge V2 \wedge V3 \wedge V4 \wedge V5 \wedge V6 .
\]
\end{theorem}

\begin{proof}
$(\Rightarrow)$ Suppose the certificate is complete. By definition
of completeness, the code is redundant (Principle~\ref{pr:redundancy}),
which requires V1. Convergence is required for the redundancy to be
meaningful, giving V2. Coherence ensures that the layers impose
compatible constraints, giving V3. Independence ensures that the
reconstruction is not an artifact of the code itself, giving V4.
Traceability and formalisation are the operational realisation of
auditability (Principle~\ref{inv:provenance}) and of \cite[Level II]{paperIII},
giving V5 and V6.

$(\Leftarrow)$ Suppose V1--V6 hold. By V1 and V2, the code is
redundant. By V3 and V4, the structure is coherent and independent.
By V5 and V6, the trace is auditable and formally verifiable.
Therefore the certificate is complete.
\end{proof}

% ------------------------------------------------------------
\section{Terminal States and the Six-Level Taxonomy}
\label{sec:terminal}
% ------------------------------------------------------------

\subsection{Mapping to Paper III}

The three terminal states of the protocol map onto the six-level
taxonomy of \cite[\S5]{paperIII}:

\begin{center}
\begin{tabular}{c|l}
  \toprule
  Protocol state & Paper III level \\
  \midrule
  \STCLOSED (V1--V6) & Level I (complete proof) \\
  \STCLOSED (V1, V3--V6) & Level II (rigorous formalisation) \\
  \STCLOSED (V1, V3--V5) & Level IV (structural characterisation) \\
  \STOPEN & Level V (numerical + limits identified) \\
  \STOBSTR & Level VI (structural obstruction) \\
  \bottomrule
\end{tabular}
\end{center}

\begin{remark}
\textsc{Closed} does not automatically mean Level~I. The
certificate's level depends on which of V1--V6 hold. Only when all
six hold is the certificate at Level~I.
\end{remark}

\subsection{Open Certificates}

An \STOPEN{} certificate is issued when V4 (independence) or V6
(formalisation) fails but V1, V2, V3, and V5 hold. The certificate
records:
\begin{itemize}
  \item the residual classes $\Ccal^{(\infty)}$;
  \item the candidate representations that might separate them;
  \item the specific missing condition (V4 or V6);
  \item the estimated effort or technique required.
\end{itemize}

\subsection{Obstructed Certificates}

An \STOBSTR{} certificate is issued when a Type~III rupture
(\cite[Definition 3.6]{paperII}) is identified. The certificate
records:
\begin{itemize}
  \item the layer $\lambda$ at which the rupture occurs;
  \item the criterion $j \in \Jset$;
  \item the residual class $\Ccal^{(\infty)}$;
  \item the reason no admissible extension of the representation
        family separates the class.
\end{itemize}

% ============================================================
\part{The Certification Document}
\label{part:certificate}
% ============================================================

\section{Specification}
\label{sec:certificate}

The output of the protocol is not a paper. It is a structured
document called a \emph{certification document}.

\begin{definition}[Certification document]\label{def:cert}
A certification document is a tuple
\[
  \mathfrak{C} = (\mathrm{Code}, \Mcal_{\rm imp}, \Ob_{\rm rec},
    \Rcal_\infty, \Acal, \Vmat, \Tcal, \mathrm{state}),
\]
where:
\begin{itemize}
  \item $\mathrm{Code}$ is the source code, unmodified;
  \item $\Mcal_{\rm imp}$ is the inferred model;
  \item $\Ob_{\rm rec}$ is the reconstructed continuous object;
  \item $\Rcal_\infty$ is the final representation family;
  \item $\Acal$ is the audit matrix;
  \item $\Vmat$ is the validation matrix;
  \item $\Tcal$ is the provenance trace;
  \item $\mathrm{state} \in \{\STCLOSED, \STOPEN, \STOBSTR\}$ is
        the terminal state.
\end{itemize}
\end{definition}

\subsection{Required Sections}

The document contains the following sections, in order:

\begin{enumerate}[label=\textbf{\S\arabic*.}]
  \item \textbf{Code.} The source code, unmodified, with line
        numbering for reference.
  \item \textbf{Inferred model.} The equation, domain, boundary
        conditions, and hypotheses, with justification.
  \item \textbf{Reconstructed object.} The continuous object
        $\Ob_{\rm rec}$ with operator, functional space, and
        conditions.
  \item \textbf{Layer analysis.} One subsection per layer: spectral,
        geometric, topological, algebraic, variational, numerical,
        analytic.
  \item \textbf{Independent re-discretisation.} The alternative
        scheme, its comparison to the original, convergence rate.
  \item \textbf{Audit matrix.} The matrix $\Acal$ with states.
  \item \textbf{Validation matrix.} The matrix $\Vmat$ with
        epistemic states.
  \item \textbf{Provenance trace.} The directed graph
        $\Tcal$ with edges from each result to its origin.
  \item \textbf{Terminal certificate.} The state declaration
        with justification and level in the taxonomy of
        \cite{paperIII}.
\end{enumerate}

\subsection{Self-Containment}

\begin{principle}[Certificate self-containment]\label{pr:self-contained}
The certification document is \emph{self-contained} if a third
party can reproduce every step of the analysis from the document
alone, without access to the original author, the original
execution, or any additional code beyond $\mathrm{Code}$.
\end{principle}

Self-containment is the operational criterion of completeness of
the certificate. It fails if any step requires information not
present in the document.

% ============================================================
\part{Worked Cases}
\label{part:cases}
% ============================================================

\section{Case 1: Finite Differences for the Dirichlet Laplacian}
\label{sec:case-closed}

\subsection{The Code}

A code fragment implementing the finite-difference Laplacian on
$(0,1)$ with Dirichlet boundary conditions:

\begin{lstlisting}[language=Python]
import numpy as np

def laplacian(N):
    h = 1.0 / (N + 1)
    diag = 2.0 * np.ones(N) / h**2
    off = -1.0 * np.ones(N - 1) / h**2
    L = np.diag(diag) + np.diag(off, 1) + np.diag(off, -1)
    return L

def eigenvalues(N):
    L = laplacian(N)
    return np.linalg.eigvalsh(L)
\end{lstlisting}

The output $\mathrm{out}$ is the sorted array of eigenvalues for
$N = 32, 64, 128, 256$.

\subsection{C0: Code Reading}

Governing hypothesis: the code approximates the operator
$-d^2/dx^2$ on $(0,1)$ with Dirichlet boundary conditions. Implicit
choices: central differences of order 2; uniform grid; eigenvalues
computed via \texttt{eigvalsh}.

\subsection{C1: Model Inference}

$\Mcal_{\rm imp}$: $u \mapsto -u''$ on $H^2(0,1) \cap H_0^1(0,1)$.

\subsection{C2: Classical Theory Identification}

Status: \textsc{Known}. The operator is the Dirichlet Laplacian;
its spectral theory is classical.

\subsection{C3: Structural Layer Construction}

\begin{itemize}
  \item Spectral: $\lambda_n = n^2 \pi^2$, $\phi_n(x) = \sin(n \pi x)$.
  \item Geometric: $M = [0,1]$, $g = dx^2$, $\Delta_g = d^2/dx^2$.
  \item Topological: $H_0 = \Z$, $H_1 = 0$.
  \item Algebraic: $L = D^\ast D$, hence $L \succeq 0$.
  \item Variational: $E[u] = \frac{1}{2}\int u_x^2$, so
        $L = -\delta E/\delta u$.
\end{itemize}

\subsection{C4: Independent Re-discretisation}

Scheme: spectral method with Chebyshev polynomials. Outputs agree
with the finite-difference outputs to $\mathcal{O}(h^2)$.

\subsection{C5: Cross-Audit}

All pairs pass: $a_{ij} = \STPASS$ for all $(i,j)$.

\subsection{C6: Code Reconstruction}

Reconstruction: the finite-difference scheme is recovered from the
continuous operator via standard second-order discretisation.
Structural and output comparisons match.

\subsection{C7: Formalization}

The definition of $L$, its self-adjointness, and its spectrum are
formalised in Lean~4 using Mathlib.

\subsection{Certificate}

Terminal state: \STCLOSED. All of V1--V6 hold. Level in the
taxonomy: Level~I.

% ------------------------------------------------------------
\section{Case 2: An Obstructed Code}
\label{sec:case-obstr}

\subsection{The Code}

A code fragment implementing a diffusion equation with a spurious
advection term:
\begin{lstlisting}[language=Python]
def step(u, dt, h, kappa, epsilon):
    lap = (np.roll(u, -1) - 2*u + np.roll(u, 1)) / h**2
    adv = (np.roll(u, -1) - np.roll(u, 1)) / (2*h)
    return u + dt * (kappa * lap + epsilon * adv)
\end{lstlisting}

The code claims to solve $u_t = \kappa u_{xx}$ but includes an
advection term $\varepsilon u_x$.

\subsection{C0--C3}

C0 identifies the extra term $\varepsilon u_x$. C1 infers the
model $u_t = \kappa u_{xx} + \varepsilon u_x$. C2 classifies it as
a variant of advection-diffusion. C3 builds the layers.

The spectral layer reveals a constant shift
$-\varepsilon^2/(4\kappa)$ in the eigenvalues, incompatible with
pure diffusion. The variational layer reveals that the operator is
not self-adjoint when $\varepsilon \neq 0$. The geometric layer
reveals a preferred direction $x \to 1-x$ that breaks the isotropy
hypothesis.

\subsection{C4--C5}

C4 builds an independent scheme. C5's audit matrix shows a
\STFAIL{} at $(\Sp, \Al)$: the spectral prediction differs from
the variational prediction by a term proportional to
$\varepsilon^2$.

\subsection{C6: Classification}

The discrepancy is Type~III: no admissible extension of the
representation family can identify the object, because the isotropy
hypothesis of the problem is incompatible with $\varepsilon \neq 0$.

\subsection{Certificate}

Terminal state: \STOBSTR. The obstruction is localised at the pair
$(\Sp, \Al)$ under the symmetry constraint of the geometric layer.
Level in the taxonomy: Level~VI.

% ============================================================
\part{Backend Architecture}
\label{part:backend}
% ============================================================

\section{Division of Labor}
\label{sec:backend-division}

The protocol is implementation-independent. One instantiating
backend is a hybrid Python--Rust system.

\begin{itemize}
  \item \textbf{Python}: mathematical instrumentation. Symbolic
        (SymPy), numeric (NumPy, SciPy), plotting, external
        libraries.
  \item \textbf{Rust}: structural audit. Constraints, coherence
        checks, provenance, certificate emission.
\end{itemize}

\subsection{Boundary}

The boundary is a typed structure
\[
  \mathrm{RepresentationRecord} = (\text{kind}, \text{data},
    \text{assumptions}, \text{provenance}, \text{state}).
\]
Python produces records; Rust consumes, checks, and emits audit
results.

\subsection{Cycle Mapping}

\begin{center}
\begin{tabular}{c l l}
  \toprule
  Cycle & Python role & Rust role \\
  \midrule
  C0 & parse code & validate schema \\
  C1 & infer model & store inferred model \\
  C2 & search literature & store classification \\
  C3 & compute layers & store representation family \\
  C4 & re-discretise & compare records \\
  C5 & --- & build audit matrix \\
  C6 & --- & attempt reconstruction \\
  C7 & formalise (Lean) & manage proofs \\
  \bottomrule
\end{tabular}
\end{center}

\subsection{Directory Layout}

\begin{lstlisting}[language=bash]
protocol/
  core/                  # Rust
    code_layer.rs
    inference.rs
    constraint.rs
    coherence.rs
    classification.rs
    reconstruction.rs
    certificate.rs
    provenance.rs
  python/
    c0_reading.py
    c1_inference.py
    c3_spectral.py
    c3_variational.py
    c3_geometric.py
    c4_rediscretization.py
  problems/
    laplacian.py
    advection_diffusion.py
  certificates/
    closed/
    open/
    obstructed/
\end{lstlisting}

% ============================================================
\part{Closing}
\label{part:closing}
% ============================================================

\section{Summary}

The protocol presented in this paper is an operational method for
analysing mathematical problems whose point of entry is executable
code. It:
\begin{itemize}
  \item treats the code as an artifact, not as a mathematical
        object;
  \item reverse-engineers the continuous object that the code
        approximates;
  \item validates the approximation via an independent
        re-discretisation;
  \item certifies the analysis via six validity conditions
        V1--V6;
  \item produces a \emph{certification document} rather than a
        paper;
  \item terminates in \STCLOSED, \STOPEN, or \STOBSTR.
\end{itemize}

The central criterion is \emph{code redundancy}: the code is fully
analysed when its outputs are predictable from the reconstructed
structure without execution.

\section{Position Relative to Companion Papers}

\begin{center}
\begin{tabular}{l l}
  \toprule
  Paper & Role in the program \\
  \midrule
  \cite{framework} & Metamodel, layers, transitions, reconstruction. \\
  \cite{paperI}    & Constraining, spectral emergence, trichotomy. \\
  \cite{paperII}   & Separation, ruptures, coherence, expansion. \\
  \cite{paperIII}  & Auditability, wear, taxonomy, trace. \\
  This paper       & Operational protocol for code-first analysis. \\
  \bottomrule
\end{tabular}
\end{center}

The five papers form a coherent program:
\begin{itemize}
  \item the framework paper defines the mathematical vocabulary;
  \item Paper~I is the operational core;
  \item Paper~II is the meta-theoretic counterpart (limits);
  \item Paper~III is the epistemological foundation;
  \item this paper is the code-first operational realisation.
\end{itemize}

\section{Closing Remark}

The protocol does not replace proof, formalisation, or peer review.
It organises the interaction between an artifact and its
mathematical interpretation so that the interpretation is explicit,
auditable, and---when possible---mechanically verifiable. The
central slogan, in one line:

\begin{quote}
\emph{The code is the question, not the answer. The protocol is the
question's formal reconstruction.}
\end{quote}

% ------------------------------------------------------------
\appendix
% ------------------------------------------------------------

\section{Notation Summary}
\label{app:notation}

\begin{center}
\begin{tabular}{c l l}
  \toprule
  Symbol & Meaning & Reference \\
  \midrule
  $\Ob$ & Master object & \cite[Def.~1.1]{framework} \\
  $\Lam$ & Layer set & \cite[Def.~1.2]{framework} \\
  $\Lam^+$ & Extended layer set & Def.~\ref{def:extended-layers} \\
  $\Cd$ & Code layer & Def.~\ref{def:extended-layers} \\
  $\Acal$ & Computational artifact & Def.~\ref{def:artifact} \\
  $\Ocal_\Cd$ & Code datum & Def.~\ref{def:code-datum} \\
  $\iota_{\Cd \to \lambda}$ & Inference morphism & Def.~\ref{def:inference} \\
  $\Mcal_{\rm imp}$ & Inferred model & Def.~\ref{def:inferred-model} \\
  $\mathrm{Ep}$ & Extended epistemic states & Def.~\ref{def:epistemic-states} \\
  $\mathfrak{C}$ & Certification document & Def.~\ref{def:cert} \\
  V1--V6 & Validity conditions & \S\ref{sec:validity} \\
  C0--C7 & Protocol cycles & \S\ref{sec:protocol-overview} \\
  \bottomrule
\end{tabular}
\end{center}

\section{Bibliography}
\label{app:biblio}

\begin{thebibliography}{99}

\bibitem{framework}
[Framework paper],
\emph{A Multirepresentational Construction--Reconstruction
Methodology for Mathematical Structures}, Parts I--V, 2024.

\bibitem{paperI}
[Paper I],
\emph{Structural Constraining and Spectral Emergence: A
Multirepresentational Methodology for Nonlinear Evolution
Equations}, 2024.

\bibitem{paperII}
[Paper II],
\emph{On the Limits of Structural Constraining: Separation,
Coherence, and the Emergence of New Representations}, 2024.

\bibitem{paperIII}
[Paper III],
\emph{Structural Auditability in Mathematical Construction: A
Methodological Position}, 2024.

\bibitem{reed-simon}
M.~Reed, B.~Simon,
\emph{Methods of Modern Mathematical Physics, Vol.~I: Functional
Analysis}, Academic Press, 1980.

\bibitem{engl}
H.~W.~Engl, M.~Hanke, A.~Neubauer,
\emph{Regularization of Inverse Problems}, Kluwer, 1996.

\bibitem{hairer}
E.~Hairer, C.~Lubich, G.~Wanner,
\emph{Geometric Numerical Integration}, Springer, 2006.

\bibitem{lax-rich}
P.~D.~Lax, R.~D.~Richtmyer,
Survey of the stability of linear finite difference equations,
\emph{Comm.~Pure Appl.~Math.}~9 (1956), 267--293.

\bibitem{mathlib}
The Mathlib Community,
\emph{Mathlib: A unified library of mathematics formalized in Lean},
2024.

\end{thebibliography}

\end{document}